\subsection*{Groups, actions and proof inputs}
Throughout, $F=\F$ is the free profinite group on $x,y$, with
$z=(xy)^{-1}$. We use the completed free product
$\B=C_2\amalg C_3=\langle s,r\mid s^2=r^3=1\rangle$,
$\tau=sr$, and the open level-two free subgroup with
$x=(sr)^{-2}$, $y=sxs$. The symbol $F'$ means the closed derived
subgroup. For a profinite group $K$, $R_{\rm nil}(K)$ is the
intersection of the kernels of all its continuous homomorphisms
to finite nilpotent groups; $O^p(K)$ uses finite $p$-groups.
All group homomorphisms and actions below are continuous unless
explicitly described as discrete.

The group $\GZ$ consists of normalized pairs $(\lambda,f)$,
$\lambda\in\widehat\Z^\times$, $f\in F'$, whose actual
free-group map is an automorphism and which satisfy the two-cycle
and three-cycle equations. The corresponding outer map is represented by
$x\mapsto x^\lambda$, $y\mapsto f^{-1}y^\lambda f$.
We write $\GT$ for the subgroup satisfying the pentagon and
$\HH=\ker(\lambda:\GZ\to\widehat\Z^\times)$.
The embedded absolute Galois group is $G_{\Q}$, and
$G_{\Q^{\rm ab}}=G_{\Q}\cap\HH$. For general multipliers a
formula on $\HH$ is extended only after its ambiguity is proved
independent of the chosen Galois lift of the multiplier.

For $h\in\HH$ we use the normalized extension fixing $\tau$.
Its positive pure frames satisfy
\[
 h(s)=\beta_hs\beta_h^{-1},\qquad
 h(r)=A_hrA_h^{-1},\qquad A_h,\beta_h\in F.
\]
The factor-centralizer theorem of
\citet[Theorem B$'$]{HerfortRibes1985} gives uniqueness in this
normalization. The quotient $q_{12}$ has target
$Q_{12}=\SL_2(\Z/12)/\{\pm I\}$, of order 576.
The corrected signed coordinates of the supplied R4 take values
in $C_6\times C_4$, and their parity characters are
$a_3=\epsilon_3T$ and $a_2=\epsilon_2W$. These are characters
of two different frame coordinates; they are not matrix entries
of one scalar frame. The reference models are
$E_3:y^2=x^3+1$, $E_4:y^2=x^3-x$, and
$r_0=(3/4)^{1/4}>0$.

We use the following explicitly stipulated, internally refereed
inputs from the supplied development \citep{SuppliedReferees2026}.
Their names are fixed throughout, so that the dependency register
can be checked without treating an unrefereed row as a theorem.
\begin{description}
\item[R1.] LM7\_PUSH D.50 with D.44: the peripheral lifting theorem
from the absolute relative-pronilpotent modular or free source;
in particular the actual multiplier-one pro-$\ell$ pair lifts
for $\ell\ge5$, with its complete pro-$\ell$ image retained.
\item[R2.] LM7\_PUSH E.58 and its referee report: the corresponding
relative lifting with specified conjugators for every open normal
subgroup. For $U=F'F^3$ it gives an element acting identically on
$\B/R_{\rm nil}(U)$ with $q_{12}(A)=[5I]$,
$q_{12}(\beta)=1$, and parities $(1,0)$.
\item[R3.] LM7\_PUSH E.61: the modulus-twelve shear covariance and
formula for $q_{12}(f)$, with $t=8\rho$ and
$U_0=\left(\begin{smallmatrix}1&1\\0&1\end{smallmatrix}\right)$.
The source's sign convention for $\rho$ is audited separately below.
\item[R4.] LM7\_PUSH E.65, with both referee repairs included:
$T,W$ are homomorphisms on $\HH$, and the sum of their parity
characters extends to the continuous character $\Delta$ on $\GZ$.
It is zero on $G_{\Q}$ and one on the R2 element. The two-model
Galois comparison and pure $r$-frame existence belong to this input;
E.62 is not a dependency.
\item[R7.] LM4 D.32 and D.34 with the two supplied referee reports:
the stated torsion assertions and the reduction of the involution
value to its multiplier-one partner.
\end{description}
Other background assertions are proved below, corrected, or retained
as exact unresolved claims. The fields and coordinates denoted by
similar letters in other returns are always given their source prefix.
In particular the finite-field statement is \emph{LM7 D.25}.

The numerical experiments in this paper use exact finite arithmetic,
not probabilistic evidence for a profinite lift. A finite solution
without a lifting theorem is called a shadow. The final unresolved
question is a statement about actual elements of $\GT$.
